% GENERATED from appendix_mcu.in.tex by tools/make_tables.py
% Appendix of the paper: microcontroller implementations on the nRF52840 and
% the nRF52832. It is placed after \appendix and before the references, and
% refers to Sections~\ref{subsec:minimality}, \ref{subsec:usecases} and
% \ref{sec:performance} of the paper. Every number in this file is filled in
% by tools/make_tables.py from results/nrf52840/ and results/nrf52832/ (fields
% in double angle brackets; the values are listed in results/values.txt), the
% tables as well.
% Bibliography keys: nrf52840cryptocell, nrf52832ecb, armCortexM4TRM and HCTR2
% are those of the paper; wolfSSLCC310, TFM-cc3xx, Unicorn, BearSSL and
% ShoupGCM are in paper/mcu_refs.bib.
% TODO(authors): once run on boards, replace the estimated cycles by the DWT
% cycle counts that the same firmware prints (see the README) and drop the
% estimation paragraphs that no longer apply.

\section{Microcontroller Implementations}
\label{app:mcu}

This appendix complements the measurements of Section~\ref{sec:performance}
with the two devices of Section~\ref{subsec:usecases}. We implemented both
modes for these two microcontrollers of the Nordic nRF52 series, which share
the same Arm Cortex-M4F core at 64~MHz, without instructions for AES, SHA-2,
SHA-3 or carry-less multiplication~\cite{armCortexM4TRM}, and differ in their
cryptographic hardware (Table~\ref{tab:mcu-devices}). The
nRF52840~\cite{nrf52840cryptocell} contains an Arm CryptoCell-310 coprocessor that
computes AES-128 in both directions and SHA-256, but has no GHASH unit (its
AES engine supports the ECB, CBC, CTR, CMAC and CCM modes, not GCM), so
POLYVAL must be computed in software: it is a device with SHA-2 and AES
hardware that lacks carry-less multiplication. The
nRF52832~\cite{nrf52832ecb} only has an AES-128 peripheral that encrypts single
blocks (ECB), and no hash hardware: it is a device with forward-only AES. In
each case, we give every mode the hardware support it can use.

\begin{table}[!hbt]
\caption{The two devices and the hardware used by each mode. Both have an Arm
Cortex-M4F core at 64~MHz without instructions for AES, SHA-2, SHA-3 or
carry-less multiplication.}
\label{tab:mcu-devices}
\centering
\scriptsize
\renewcommand{\arraystretch}{1.25}
\setlength{\tabcolsep}{3pt}
\begin{tabular}{>{\raggedright\arraybackslash}p{1.9cm} >{\raggedright\arraybackslash}p{4.75cm} >{\raggedright\arraybackslash}p{4.75cm}}
\hline
 & nRF52840 & nRF52832 \\
\hline\hline
Memory & 1~MiB flash, 256~KiB RAM & 512~KiB flash, 64~KiB RAM \\
AES hardware & CryptoCell-310: AES-128, encryption and decryption; ECB
  peripheral (encryption) & ECB peripheral: AES-128, encryption of single
  blocks only \\
Hash hardware & CryptoCell-310: SHA-1, SHA-224, SHA-256 & none \\
Carry-less multiplication & none & none \\
\hline
Concertina & AES and SHA-256 compression on the CryptoCell; HMAC-SHA3-256 with
  Keccak-$f$ in software & AES on the ECB peripheral; HMAC in software \\
HCTR2 & AES (both directions) on the CryptoCell; POLYVAL in software &
  forward AES on the ECB peripheral; POLYVAL and the inverse cipher (one
  block per decryption) in software \\
\hline
\end{tabular}
\end{table}

\paragraph{Implementation.}
All modes share one AES layer with a backend for each AES unit, a 32-bit
table implementation in software, and AES-XTS is listed as a narrow-block
reference. The CryptoCell is driven through its documented registers by a
bare-metal driver: AES-128 in ECB mode over DMA, encrypting or decrypting up
to 16 blocks per transfer, and the SHA-256 compression function, loaded with
a chaining value and fed whole blocks over DMA without hardware padding.
Neither the counter of Concertina, $E_K(S \oplus \langle i \rangle)$, nor the
XCTR counter of HCTR2, $E_K(S \oplus \mathrm{le64}(i))$, is the CTR mode of
the CryptoCell, so both modes prepare their counter blocks on the CPU and
encrypt them in batches of 16. HMAC-SHA-256 keeps its precomputed inner and
outer states and pads in software, so that each of the two HMAC evaluations
of a message costs two or three DMA transfers (the whole blocks of the
input, the padded last block or blocks, and the outer hash); the KDF of the
key setup uses the same HMAC. The ECB peripheral
processes one block per operation. HMAC-SHA3-256 uses a bit-interleaved
Keccak-$f$ in software on both chips. For HCTR2, we measure three POLYVAL
implementations: a constant-time one with $32 \times 32$-bit
multiplications on bit-masked operands, Karatsuba multiplication and a
Montgomery reduction by shifts; and the table-based method of
Shoup~\cite{ShoupGCM} with 4-bit (256~bytes per key, as in mbed~TLS) and
8-bit (4~KiB per key) tables, which index memory with secret data. On the
nRF52832, HCTR2 decryption evaluates its single inverse-cipher block either
with a table implementation or with the constant-time bitsliced
implementation of BearSSL~\cite{BearSSL}. As in
Section~\ref{sec:performance}, HCTR2 was validated against the 96 POLYVAL,
180 XCTR and 700 HCTR2 test vectors of its reference
implementation~\cite{HCTR2} and Concertina against its 90 known answers, here
in every measured configuration (the AES hardware, which supports only
128-bit keys, with the vectors for AES-128); the firmware runs
13{,}427 checks on the nRF52840 and 11{,}604 on the nRF52832 before
timing. The code was compiled with GCC~13.2 (\texttt{-O2
-fno-schedule-insns}); all measurements use AES-128 and a 16-byte tweak. As
for the distinguishers of Section~\ref{subsec:minimality}, the firmware, the
emulator described below and the measurements were developed with the
assistance of a large language model (Claude, Anthropic); they are part of
the repository referenced in Section~\ref{sec:performance}.

\paragraph{Measurement method: exact and estimated quantities.}
Without access to boards, we executed the firmware in an emulator of the two
chips: an instruction-set emulator of the Cortex-M4 (Unicorn~\cite{Unicorn},
based on QEMU) with the memory map of each chip and functional models of the
peripherals the firmware uses. The model of the CryptoCell implements the
documented register interface of its AES and HASH engines and their DMA
(register layout from Nordic's device description and product
specification; key and chaining-value word order as in Arm's drivers for the
CryptoCell-312~\cite{TFM-cc3xx}, whose registers have the same offsets); it
rejects accesses that the documentation forbids, such as DMA from flash, and
computes its results with implementations independent of the firmware. The
same firmware image is built to run unchanged on the boards, where it prints
cycle counts. We distinguish three kinds of results.
\begin{itemize}
\item \emph{Exact}, in the sense that a board running the same image gives
the same values: code size, static RAM, context size, peak stack depth
(measured by painting the stack), the number of executed instructions apart
from the iterations of the drivers' busy-wait loops (which the emulator
leaves after one test), and the number and size of the hardware operations
of each message (for data in RAM, as in the benchmark; the DMA cannot read
flash, and the driver copies such input through a 256-byte buffer). For
instance, Concertina with HMAC-SHA-256 on the nRF52840 makes 4 SHA-256 transfers
(4 blocks) and 4 AES transfers for a
32-byte message, and 6 SHA-256 transfers
(130 blocks) and 20 AES transfers
(258 blocks) for 4096 bytes.
\item \emph{Established against the models only}: functional
correctness of the hardware backends. All test vectors pass with every
backend in the emulator, but the CryptoCell driver has not run on silicon;
in particular, the order of the chaining-value words and the meaning of the
hash endianness setting, whose description in Nordic's device file
contradicts its reset value, are taken from Arm's drivers. The self-tests
that the firmware runs before timing would detect a mismatch on a board.
\item \emph{Estimated}: execution time. It is the sum of a CPU part and a
hardware part. The CPU part converts the executed instructions into cycles
with the instruction timings of the Cortex-M4 Technical Reference
Manual~\cite{armCortexM4TRM},
with a pipeline refill of $P = 1$, 2 or 3 cycles per taken branch; this
gives between 1.18 and 1.90 cycles per instruction for the
modes of Table~\ref{tab:mcu-speed}, and does not model flash wait states. The
hardware part charges each hardware operation a latency, while the CPU waits
for it (Table~\ref{tab:mcu-latency}): for the ECB peripheral, the typical
6~$\mu$s per block of the product specification~\cite{nrf52832ecb}; for the CryptoCell, whose
latencies are not specified, a per-block cost and a per-transfer cost in
three profiles. The \emph{fast} profile takes the 4 cycles per AES round
that the AES engine reports in its configuration
register~\cite{nrf52840cryptocell} and one cycle per SHA-256 round; the
\emph{central} profile takes the AES-128-CBC and SHA-256 throughputs
measured by wolfSSL with the CryptoCell of an nRF52840
board~\cite{wolfSSLCC310}, 17.3 and 25.9~MB/s with 1~MB = $2^{20}$ bytes
(the unit of wolfCrypt's benchmark), that is 56 cycles per AES block and
151 cycles per SHA-256 block, which include the overhead of their
software; the \emph{slow} profile doubles the central one. The cost of a
transfer (start, completion, pipeline fill) is not documented, and we take
16, 64 and 256 cycles.
\end{itemize}
Table~\ref{tab:mcu-speed} uses $P = 2$ and the central profile. Since the
CPU part does not depend on the latencies, each time is a linear function of
them, and the counts of hardware operations give the time under any other
latency; Table~\ref{tab:mcu-robust} lists the extremes over all combinations.
The cost per transfer covers only the hardware: the driver code around
each transfer is executed and counted in the CPU part, and a vendor library
with locking or power management per call would add to it; we therefore also
give the break-even transfer cost.

\begin{table}[!hbt]
\caption{Latencies charged for the hardware operations, in cycles at 64~MHz.
The CryptoCell profiles are estimates (see text); the central profile is
used unless stated otherwise.}
\label{tab:mcu-latency}
\centering
\small
\renewcommand{\arraystretch}{1.25}
\begin{tabular}{lrrr}
\hline
 & Fast & Central & Slow \\
\hline\hline
CryptoCell: per DMA transfer (AES or SHA-256) & 16 & 64 & 256 \\
CryptoCell: per 16-byte AES-128 block & 40 & 56 & 112 \\
CryptoCell: per 64-byte SHA-256 block & 64 & 151 & 302 \\
\hline
ECB peripheral (nRF52832): per block & 192 & 384 & 768 \\
\hline
\end{tabular}
\end{table}

\begin{table}[!hbt]
\caption{Encryption cost on the nRF52840 and the nRF52832 (Arm Cortex-M4F, 64~MHz) with AES-128 and a 16-byte tweak, in estimated cycles per byte: the CPU cycles of the executed instructions under the Cortex-M4 timing model ($P = 2$), plus the latency of each hardware operation (CryptoCell-310: central profile of Table~\ref{tab:mcu-latency}; ECB peripheral: typical $t_\mathrm{ECB}$ of 6~$\mu$s). The software rows are the same on both chips.}
\label{tab:mcu-speed}
\centering
\small
\renewcommand{\arraystretch}{1.25}
\setlength{\tabcolsep}{4pt}
\begin{tabular}{lrrrrrr}
\hline
 & \multicolumn{6}{c}{Message length (bytes)} \\
Mode & 32 & 64 & 256 & 512 & 1024 & 4096 \\
\hline\hline
\multicolumn{7}{l}{\emph{nRF52840: AES-128 and SHA-256 on the CryptoCell-310}} \\
AES-XTS$^\dagger$                                    & 172 & 91.3 & 30.9 & 21.3 & 16.6 & 13.0 \\
HCTR2, constant-time POLYVAL                         & 197 & 155 & 124 & 119 & 117 & 115 \\
HCTR2, 4-bit POLYVAL tables                          & 220 & 184 & 157 & 153 & 151 & 150 \\
HCTR2, 8-bit POLYVAL tables                          & 159 & 113 & 78.1 & 72.9 & 70.2 & 68.3 \\
Concertina (HMAC-SHA-256)                            & 272 & 146 & 48.3 & 30.8 & 22.1 & 15.5 \\
Concertina (HMAC-SHA3-256)                           & 2077 & 1048 & 383 & 310 & 273 & 239 \\
\hline
\multicolumn{7}{l}{\emph{nRF52832: forward AES-128 on the ECB peripheral}} \\
AES-XTS$^\dagger$                                    & 61.9 & 48.6 & 38.6 & 37.0 & 36.1 & 35.5 \\
HCTR2, constant-time POLYVAL                         & 149 & 145 & 142 & 142 & 141 & 141 \\
HCTR2, 4-bit POLYVAL tables                          & 172 & 174 & 175 & 176 & 176 & 176 \\
HCTR2, 8-bit POLYVAL tables                          & 111 & 103 & 96.2 & 95.2 & 94.7 & 94.3 \\
Concertina (HMAC-SHA-256)                            & 541 & 290 & 166 & 145 & 134 & 127 \\
Concertina (HMAC-SHA3-256)                           & 2042 & 1041 & 401 & 332 & 298 & 265 \\
\hline
\multicolumn{7}{l}{\emph{Software only (either chip)}} \\
AES-XTS$^\dagger$                                    & 95.1 & 76.3 & 62.1 & 59.8 & 58.6 & 57.7 \\
HCTR2, constant-time POLYVAL                         & 166 & 163 & 160 & 160 & 160 & 160 \\
HCTR2, 8-bit POLYVAL tables                          & 128 & 120 & 115 & 114 & 113 & 113 \\
Concertina (HMAC-SHA-256)                            & 582 & 320 & 187 & 165 & 154 & 145 \\
Concertina (HMAC-SHA3-256)                           & 2082 & 1070 & 423 & 352 & 317 & 284 \\
\hline
\multicolumn{7}{l}{\footnotesize $^\dagger$Narrow-block mode, listed as a reference only.}
\end{tabular}
\end{table}

\paragraph{nRF52840: AES and SHA-256 in hardware.}
With both AES and SHA-256 on the CryptoCell, Concertina with HMAC-SHA-256
needs 15.5 cycles per byte for 4096-byte messages, close to the
13.0 cycles per byte of AES-XTS. HCTR2 also computes its AES calls
on the CryptoCell, but POLYVAL remains on the CPU: one block costs
853 cycles with the constant-time implementation and
476 with 8-bit tables, and POLYVAL accounts for
92\% and 87\% of HCTR2 at 4096
bytes. HCTR2 needs 115, 150 and 68.3
cycles per byte with the constant-time, 4-bit and 8-bit POLYVAL, so that
Concertina is 7.4, 9.6 and
4.4 times as fast at 4096 bytes. The 4-bit tables are
slower than the constant-time multiplication on this core, whose
$32 \times 32$-bit multiplier takes a single cycle. Since the CPU
part of HCTR2 alone amounts to 111 cycles per byte with the
constant-time POLYVAL and 64.5 with 8-bit tables, no faster
AES hardware would change this. Conversely, for HCTR2 to match
Concertina, POLYVAL would have to cost less than about 55
cycles per 16-byte block at 4096 bytes and 230 at 256 bytes
(98 and 351 under the least favorable of
our timing assumptions), against 476 for our fastest
implementation; at 4096 bytes, this is fewer cycles than the 64 word loads
and 64 XOR operations of the table lookups alone in Shoup's method with 8-bit
tables. For short messages, the fixed cost of
Concertina dominates: its two HMAC evaluations take 4
SHA-256 transfers and its four Feistel rounds 4
single-block AES transfers, each with its driver code on the CPU, and at 32
bytes Concertina needs 272 cycles per byte against 197
and 159 for HCTR2 with the constant-time and 8-bit POLYVAL.
Concertina is faster than HCTR2 with the constant-time and the 4-bit POLYVAL
from 64 bytes on, and than HCTR2 with 8-bit tables from
256 bytes on.
HMAC-SHA3-256, without SHA-3 hardware, remains slower than HCTR2 at every
length (239 cycles per byte at 4096 bytes).

\begin{table}[!hbt]
\caption{Speed-up of Concertina (HMAC-SHA-256) over HCTR2 (time of HCTR2 divided by that of Concertina; encryption, and decryption in the last row) under all combinations of the CPU timing model ($P = 1, 2, 3$) and the hardware latencies: on the nRF52840 the three CryptoCell profiles of Table~\ref{tab:mcu-latency}, on the nRF52832 $t_\mathrm{ECB}$ = 3, 6 and 12~$\mu$s. Minimum and maximum, rounded outwards; values above 1 favor Concertina.}
\label{tab:mcu-robust}
\centering
\scriptsize
\renewcommand{\arraystretch}{1.25}
\setlength{\tabcolsep}{3pt}
\begin{tabular}{lrrrrrr}
\hline
 & \multicolumn{6}{c}{Message length (bytes)} \\
POLYVAL of HCTR2 & 32 & 64 & 256 & 512 & 1024 & 4096 \\
\hline\hline
\multicolumn{7}{l}{\emph{nRF52840 (CryptoCell-310)}} \\
Constant-time                                & 0.60--0.81 & 0.85--1.22 & 1.87--3.15 & 2.70--4.92 & 3.54--7.00 & 4.68--10.38 \\
4-bit tables                                 & 0.67--0.90 & 1.01--1.43 & 2.40--3.92 & 3.52--6.20 & 4.65--8.86 & 6.09--13.20 \\
8-bit tables                                 & 0.50--0.64 & 0.64--0.86 & 1.24--1.91 & 1.74--2.90 & 2.23--4.05 & 2.86--5.94 \\
\hline
\multicolumn{7}{l}{\emph{nRF52832 (ECB peripheral)}} \\
Constant-time                                & 0.26--0.30 & 0.47--0.53 & 0.83--0.88 & 0.95--1.01 & 1.02--1.09 & 1.07--1.16 \\
4-bit tables                                 & 0.30--0.34 & 0.59--0.62 & 1.03--1.08 & 1.16--1.24 & 1.25--1.35 & 1.31--1.44 \\
8-bit tables                                 & 0.19--0.23 & 0.33--0.39 & 0.55--0.63 & 0.62--0.70 & 0.67--0.75 & 0.71--0.79 \\
Constant-time, with $\mathrm{AES}^{-1}$ (dec.) & 0.71--0.80 & 0.89--1.01 & 1.00--1.10 & 1.04--1.13 & 1.07--1.16 & 1.09--1.17 \\
\hline
\end{tabular}
\end{table}

At the tabulated lengths, these conclusions do not depend on the timing
assumptions (Table~\ref{tab:mcu-robust}), except at 64 bytes, where the
order of Concertina and HCTR2 with the constant-time POLYVAL depends on
them. From 256 bytes on, Concertina is faster than every HCTR2 variant under
every combination of the CPU model and the CryptoCell profiles, and HCTR2 is
faster at 32 bytes under every combination; at 4096 bytes, Concertina is
faster by a factor of 2.86 to
5.94 against the 8-bit tables and of
4.68 to 10.38 against the constant-time
POLYVAL. Equivalently, with the other latencies central,
Concertina remains faster than HCTR2 with 8-bit tables at 4096 bytes as long
as a SHA-256 block costs the hardware less than 1{,}813
cycles, and faster than HCTR2 with the constant-time POLYVAL as long as it
costs less than 3{,}287 cycles; a SHA-256 compression in
software takes 2{,}874 cycles on this core. At 256 bytes, the
break-even costs are 915 and 2{,}087
cycles per block. Concertina makes more CryptoCell transfers than HCTR2
(11 against 2 at 256 bytes), so a high
cost per transfer works against it: with the other latencies central, it
remains faster than HCTR2 with 8-bit tables as long as a transfer costs less
than 913 cycles at 256 bytes and 24{,}065
at 4096 bytes, against our assumption of 16 to
256.

\paragraph{nRF52832: forward-only AES.}
On the nRF52832, the ECB peripheral computes the AES calls of both modes,
and all hashing is done in software. The cost per byte then depends mostly on
the hash layer: at 4096 bytes, Concertina with HMAC-SHA-256 needs
127 cycles per byte and HCTR2 141, 176
and 94.3 with the constant-time, 4-bit and 8-bit POLYVAL.
Concertina is faster than HCTR2 with the constant-time POLYVAL from
1024 bytes on and with 4-bit tables from 256 bytes
on, by factors of 1.12 and 1.39 at
4096 bytes,
but slower than HCTR2 with 8-bit tables at every length (a ratio of
0.75 at 4096 bytes), and several times slower for short
messages. HCTR2 decryption needs one block of the inverse cipher, which the
peripheral cannot compute (Table~\ref{tab:mcu-dec}): with a table
implementation, decryption costs up to 6.6\% more than
encryption; with the constant-time implementation, which a constant-time
HCTR2 requires, it costs 8{,}131 cycles per message more than with
the table implementation, and 177\% more than encryption at 32
bytes. Against this fully constant-time HCTR2, Concertina decrypts faster
from 256 bytes on. With HMAC-SHA3-256, Concertina is slower
than every HCTR2 variant at every length on this chip as well. On this
device, Concertina is not the faster mode; it avoids the inverse cipher and
secret-indexed tables.

\begin{table}[!hbt]
\caption{Decryption cost on the nRF52840 and the nRF52832 (Arm Cortex-M4F, 64~MHz) with AES-128 and a 16-byte tweak, in estimated cycles per byte: the CPU cycles of the executed instructions under the Cortex-M4 timing model ($P = 2$), plus the latency of each hardware operation (CryptoCell-310: central profile of Table~\ref{tab:mcu-latency}; ECB peripheral: typical $t_\mathrm{ECB}$ of 6~$\mu$s). The software rows are the same on both chips.}
\label{tab:mcu-dec}
\centering
\small
\renewcommand{\arraystretch}{1.25}
\setlength{\tabcolsep}{4pt}
\begin{tabular}{lrrrrrr}
\hline
 & \multicolumn{6}{c}{Message length (bytes)} \\
Mode & 32 & 64 & 256 & 512 & 1024 & 4096 \\
\hline\hline
\multicolumn{7}{l}{\emph{nRF52840: AES-128 and SHA-256 on the CryptoCell-310}} \\
AES-XTS$^\dagger$                                    & 172 & 91.3 & 30.9 & 21.3 & 16.6 & 13.0 \\
HCTR2, constant-time POLYVAL                         & 198 & 156 & 124 & 119 & 117 & 115 \\
HCTR2, 4-bit POLYVAL tables                          & 221 & 185 & 157 & 153 & 151 & 150 \\
HCTR2, 8-bit POLYVAL tables                          & 160 & 113 & 78.2 & 72.9 & 70.2 & 68.3 \\
Concertina (HMAC-SHA-256)                            & 273 & 146 & 48.3 & 30.8 & 22.1 & 15.5 \\
Concertina (HMAC-SHA3-256)                           & 2077 & 1049 & 383 & 310 & 273 & 239 \\
\hline
\multicolumn{7}{l}{\emph{nRF52832: forward AES-128 on the ECB peripheral}} \\
AES-XTS$^\dagger$                                    & 81.4 & 68.1 & 58.1 & 56.5 & 55.6 & 55.0 \\
HCTR2, table $\mathrm{AES}^{-1}$$^\ddagger$          & 159 & 150 & 143 & 142 & 142 & 141 \\
HCTR2, constant-time $\mathrm{AES}^{-1}$$^\ddagger$  & 413 & 277 & 175 & 158 & 150 & 143 \\
HCTR2, 4-bit POLYVAL tables                          & 182 & 179 & 177 & 176 & 176 & 176 \\
HCTR2, 8-bit POLYVAL tables                          & 121 & 108 & 97.5 & 95.8 & 95.0 & 94.3 \\
Concertina (HMAC-SHA-256)                            & 542 & 290 & 166 & 145 & 134 & 127 \\
Concertina (HMAC-SHA3-256)                           & 2042 & 1041 & 401 & 332 & 298 & 265 \\
\hline
\multicolumn{7}{l}{\emph{Software only (either chip)}} \\
AES-XTS$^\dagger$                                    & 92.5 & 73.7 & 59.5 & 57.2 & 56.0 & 55.1 \\
HCTR2, constant-time POLYVAL                         & 165 & 162 & 160 & 160 & 160 & 160 \\
HCTR2, 8-bit POLYVAL tables                          & 127 & 120 & 114 & 114 & 113 & 113 \\
Concertina (HMAC-SHA-256)                            & 582 & 320 & 187 & 165 & 154 & 145 \\
Concertina (HMAC-SHA3-256)                           & 2083 & 1070 & 423 & 352 & 317 & 284 \\
\hline
\multicolumn{7}{l}{\footnotesize $^\dagger$Narrow-block mode, listed as a reference only.} \\
\multicolumn{7}{l}{\footnotesize $^\ddagger$With the constant-time POLYVAL.}
\end{tabular}
\end{table}

\begin{table}[!hbt]
\caption{Memory footprint with AES-128, in bytes: flash for key setup, encryption and decryption (including the hardware drivers); static RAM (tables, DMA buffers); per-key context; peak stack depth; and their sum (RAM). Each configuration is built with only the backends it uses; on the nRF52832, the tables of HCTR2's table-based inverse cipher are kept in flash (one block per decryption), those of AES-XTS in RAM. Key setup in $\mu$s at 64~MHz, estimated as in Table~\ref{tab:mcu-speed}.}
\label{tab:mcu-memory}
\centering
\small
\renewcommand{\arraystretch}{1.25}
\setlength{\tabcolsep}{3.5pt}
\begin{tabular}{lrrrrrr}
\hline
Mode & Flash & Static & Context & Stack & RAM & Key setup \\
\hline\hline
\multicolumn{7}{l}{\emph{nRF52840: AES-128 and SHA-256 on the CryptoCell-310}} \\
AES-XTS$^\dagger$                                    & 1{,}940 & 516 & 72 & 448 & 1{,}036 & 3.77 \\
HCTR2, constant-time POLYVAL                         & 5{,}308 & 516 & 260 & 472 & 1{,}248 & 49.8 \\
HCTR2, 4-bit POLYVAL tables                          & 2{,}864 & 516 & 372 & 320 & 1{,}208 & 64.1 \\
HCTR2, 8-bit POLYVAL tables                          & 3{,}360 & 516 & 4{,}212 & 320 & 5{,}048 & 218 \\
Concertina (HMAC-SHA-256)                            & 4{,}200 & 516 & 192 & 904 & 1{,}612 & 304 \\
Concertina (HMAC-SHA3-256)                           & 7{,}656 & 516 & 536 & 1{,}908 & 2{,}960 & 3{,}998 \\
\hline
\multicolumn{7}{l}{\emph{nRF52832: forward AES-128 on the ECB peripheral}} \\
AES-XTS$^\dagger$                                    & 2{,}268 & 1{,}592 & 232 & 496 & 2{,}320 & 76.5 \\
HCTR2, table $\mathrm{AES}^{-1}$$^\ddagger$          & 7{,}100 & 48 & 420 & 472 & 940 & 121 \\
HCTR2, constant-time $\mathrm{AES}^{-1}$$^\ddagger$  & 7{,}112 & 48 & 596 & 1{,}000 & 1{,}644 & 236 \\
HCTR2, 4-bit POLYVAL tables                          & 4{,}656 & 48 & 532 & 336 & 916 & 135 \\
HCTR2, 8-bit POLYVAL tables                          & 5{,}148 & 48 & 4{,}372 & 336 & 4{,}756 & 288 \\
Concertina (HMAC-SHA-256)                            & 4{,}448 & 48 & 192 & 1{,}308 & 1{,}548 & 901 \\
Concertina (HMAC-SHA3-256)                           & 7{,}224 & 48 & 536 & 1{,}908 & 2{,}492 & 3{,}998 \\
\hline
\multicolumn{7}{l}{\emph{Software only (nRF52840 build)}} \\
AES-XTS$^\dagger$                                    & 2{,}716 & 2{,}312 & 584 & 320 & 3{,}216 & 111 \\
HCTR2, constant-time POLYVAL                         & 6{,}636 & 2{,}312 & 596 & 472 & 3{,}380 & 158 \\
HCTR2, 8-bit POLYVAL tables                          & 4{,}684 & 2{,}312 & 4{,}548 & 400 & 7{,}260 & 326 \\
Concertina (HMAC-SHA-256)                            & 5{,}324 & 1{,}028 & 1{,}072 & 1{,}308 & 3{,}408 & 980 \\
Concertina (HMAC-SHA3-256)                           & 8{,}096 & 1{,}028 & 1{,}416 & 1{,}908 & 4{,}352 & 4{,}077 \\
\hline
\multicolumn{7}{l}{\footnotesize $^\dagger$Narrow-block mode, listed as a reference only. $^\ddagger$With the constant-time POLYVAL.}
\end{tabular}
\end{table}

\paragraph{Memory.}
Table~\ref{tab:mcu-memory} lists the memory footprint of each
configuration, built with only the backends it uses; RAM is the sum of the
static data, the context and the peak stack depth. With the CryptoCell, the
footprints of the two modes are of the same order. Concertina with
HMAC-SHA-256 needs 4{,}200 bytes of flash and
1{,}612 bytes of RAM, against
5{,}308 and 1{,}248 for HCTR2 with the
constant-time POLYVAL, whose unrolled multiplication dominates its code
size, 2{,}864 and 1{,}208 with 4-bit
tables, and 3{,}360 and 5{,}048 with
8-bit tables, which add 4~KiB to every key context. Concertina needs more
stack (904 bytes, for the HMAC computations)
and a smaller context (192 bytes). On the
nRF52832, Concertina needs no AES software at all, since it never uses the
inverse cipher, and stores its working keys unexpanded: with HMAC-SHA-256 it
takes 4{,}448 bytes of flash and
1{,}548 bytes of RAM, of which
1{,}308 are stack, mostly for the software
SHA-256. HCTR2 must keep a software inverse cipher for decryption. With
its tables in flash, since it is used for one block per decryption, HCTR2
needs 7{,}100 bytes of flash and 940
bytes of RAM with the constant-time POLYVAL, 4{,}656 and
916 with 4-bit tables, and 5{,}148
and 4{,}756 with 8-bit tables; with the constant-time
inverse cipher, it needs 7{,}112 and
1{,}644 bytes. Concertina thus needs the least flash on
this chip, more RAM than HCTR2 with a table-based inverse cipher, and less
than the fully constant-time HCTR2 or HCTR2 with 8-bit tables. Key setup takes
304~$\mu$s for Concertina with
HMAC-SHA-256 on the CryptoCell and 901~$\mu$s
on the nRF52832, where the KDF runs in software, against
49.8 and 121~$\mu$s for HCTR2
with the constant-time POLYVAL; with HMAC-SHA3-256 it takes about
4.0~ms on either chip.

\paragraph{Summary.}
The nRF52840 is an existing device with SHA-2 and AES hardware but without
carry-less multiplication. On it, Concertina instantiated with HMAC-SHA-256
outperforms HCTR2 for messages of 256 bytes and more, at the tabulated
lengths and under every timing assumption we considered, both against a
constant-time and against table-based POLYVAL implementations: at 4096
bytes by a factor of 4.4 against 8-bit tables,
7.4 against the constant-time POLYVAL and
9.6 against 4-bit tables with the central assumptions.
HCTR2 remains faster for 32-byte messages, and the memory footprints of the
two modes are of the same order. With HMAC-SHA3-256, for which neither chip
has hardware, Concertina is slower than HCTR2 on both chips. On the
nRF52832, which has forward-only AES but no hash hardware, Concertina with
HMAC-SHA-256 is not faster than HCTR2 with 8-bit POLYVAL tables; it is
faster than a fully constant-time HCTR2 from 1024 bytes on for
encryption and from 256 bytes on for decryption, and needs the
least flash and no inverse cipher. Code size, RAM and operation counts are
exact; the cycle counts are estimates, and the hardware backends have run
against models of the peripherals only: confirming both on boards remains
to be done.
